% ---------- Abstract ----------
\begin{abstract}
Machine learning models for system diagnostics require large volumes of kernel execution traces for training, yet collecting production traces is expensive due to runtime overhead, storage costs, and privacy concerns.
We present \textsc{SynthTrace}, a diffusion-based framework for generating synthetic kernel traces that can augment limited real data for downstream ML tasks.
\textsc{SynthTrace} models traces as multi-channel sequences (event types, timestamps, CPU affinity, thread identifiers, process metadata) using a Transformer-based denoising diffusion process, coupled with constraint-guided repair to enforce system invariants.

We evaluate synthetic trace quality on six diverse benchmarks through a downstream next-event prediction task.
Our results reveal strong workload dependence: for deterministic, compute-heavy workloads (scimark2), synthetic augmentation achieves 87.2\% F1-Macro at context length L=4096---only 2.6 percentage points below real-only baselines---demonstrating production viability.
In contrast, I/O-intensive workloads exhibit 25--36\% degradation, indicating fundamental challenges for current generative approaches.
Context length emerges as the dominant quality factor, with L=4096 providing +104\% relative improvement over L=256.
Constraint-guided repair consistently improves synthetic data quality by 13--18\% across configurations.
Ablation studies show that 2-channel models achieve 97--99\% of full 6-channel performance while requiring only half of computational resources.

\textsc{SynthTrace} enables practitioners to augment limited real traces with high-fidelity synthetic data for deterministic workloads, reducing data collection costs while providing actionable guidance on when synthetic augmentation is viable.
\end{abstract}